\secc Segment-based visualization for VLF lightning mapping

The described findings have direct implications, for example, on the way lightning discharge events should be displayed and localized. Given the typical dimensions and time scales on which discharges occur, it follows that even receivers operating on the VLF principle can provide meaningful descriptions of lightning discharge structures. At the same time, for the same reason, it is misleading to display lightning discharges on maps as points, which are intuitively interpreted as CG lightning strikes hitting the ground at approximately the position of the map mark. In reality, most of those points correspond to CC lightning discharges, given the actual incidence of their occurrence over CG lightning.

A better representation that considers this fact is attempted in Figure \ref[VLF_lightning_map], which shows a calculated approximate visualization of the lightning structure. The representation is based on VLF data obtained from only three measuring stations placed on cars, so it contains many simplifications. For example, the lightning structure is depicted in a 2D plane positioned at a fixed height above the ground. Additionally, it is important to note that lightning discharges do not radiate at branching points, but primarily at the locations of two-thirds recoil leaders. The drawn connections do not reflect this fact, as the depiction is currently an illustrative algorithm that, despite its limitations, is a closer description of the physical principle of lightning discharges than the widely used depiction of points in a map.

\medskip
\clabel[VLF_lightning_map]{Lightning map based on VLF data}
\picw=15cm \cinspic ./img/VLF_lightning_map.png
\caption/f An illustrative visualization of the structure of lightning discharges based on VLF data obtained from three measuring stations mounted on vehicles. The representation depicts the lightning structure in a 2D plane at a fixed height above the ground, capturing the approximate spatial distribution and progression of the lightning channels. The depiction emphasizes the complexity of the lightning and provides a closer representation of the physical nature of lightning discharges.
\medskip

The algorithm used\urlnote{https://github.com/ODZ-UJF-AV-CR/CRREAT_cars/blob/master/VLF_lightning/VLF_location.ipynb}
for processing and visualizing VLF recordings from three measurement vehicles
operates as follows.

\secc Signal preprocessing and segmentation

The recordings are first aligned according to absolute time marks derived from GNSS PPS triggers. 
Each raw signal $s_i(t)$ is normalized to zero mean and unit variance, and slow offsets are removed.
The prepared waveforms are divided into a number of contiguous segments $S_k$
chosen such that each segment contains comparable total energy $E_k$
determined relative to the weakest signal:
$$
E_k = \int_{S_k} |s_i(t)|^2 \, dt , \qquad
E_k \approx E_{\rm min} \,.
\eqmark
$$

\seccc Cross-correlation and time-difference of arrival

For every segment, the cross-correlation function between the signals from
stations $i$ and $j$ is computed as
$$
R_{ij}(\tau) = \int s_i(t)\,s_j(t+\tau)\,dt \,,
$$
and the time offset corresponding to the correlation maximum is
$$
\Delta t_{ij} = \mathop{\rm arg\,max}_{\tau} R_{ij}(\tau) \,.
$$

With one reference station 1, the delays reduce to
$$
\Delta t_i = t_i - t_1 \,.
$$

Physically impossible solutions are rejected using the condition
$$
|c \,\Delta t_i| < \Vert \vec r_i - \vec r_1 \Vert \,.
\eqmark
$$

\seccc Numerical localization

Assuming the radio wave propagates at the speed of light $c$, 
the distances between the source at position $\vec r=(x,y)$
and individual stations $\vec r_i$ satisfy
$$
\sqrt{(x - x_i)^2 + (y - y_i)^2} -
\sqrt{(x - x_1)^2 + (y - y_1)^2} = c\,\Delta t_i \,.
$$

The location $\hat{\vec r}$ is obtained by minimizing the sum of squared residuals,
$$
\hat{\vec r} = 
\mathop{\rm arg\,min}_{\vec r} \sum_i
\left(
\sqrt{(x - x_i)^2 + (y - y_i)^2} -
\sqrt{(x - x_1)^2 + (y - y_1)^2} -
c\,\Delta t_i
\right)^2 \,.
\eqmark
$$

In the presented simplified visualization, this full 3D solution
is projected onto a 2D plane at the expected cloud-base altitude.

\seccc Bearing-line visualization

Instead of showing localized points, each receiver station $i$ is represented
by a short bearing line indicating the direction of arrival of the VLF
signal.  
If the receiver measures two orthogonal electric-field components $E_x(t)$ and
$E_y(t)$, the instantaneous azimuth of arrival is
$$
\alpha_i = \mathop{\rm atan2}\bigl(E_y(t_i), E_x(t_i)\bigr) \,,
$$
and the corresponding unit direction vector
$$
\vec d_i = (\cos\alpha_i, \sin\alpha_i) \,.
$$

Each bearing line is then plotted parametrically as
$$
\vec r(s) = \vec r_i + s\,\vec d_i, \qquad 0 \le s \le L \,,
$$
where $\vec r_i$ is the station position and $L$ denotes the plotted segment length
(tens of kilometers in map scale). The ensemble of these line segments for a given
time interval represents the instantaneous spatial organization of the lightning current system,
rather than a set of isolated point impacts.

\seccc Graph-based structural visualization

The final visualization stage converts the geometrical information from the bearing lines
into an approximate channel-like structure. This is formulated as a graph construction
problem.

We denote by $\{\vec p_m\}$ a set of candidate spatial points extracted from the
bearing lines. In practice, $\vec p_m$ are either:
(i) intersections between two bearing lines from different stations, or
(ii) extrapolated continuation points along a bearing line with sufficient signal energy
in the corresponding segment.

Formally, an intersection point $\vec p_{ij}$ between stations $i$ and $j$ is
obtained by solving
$$
\vec r_i + s_i\,\vec d_i = \vec r_j + s_j\,\vec d_j \,,
$$
for the parameters $s_i, s_j$.  
If a valid pair $(s_i,s_j)$ exists with $0 \le s_i,s_j \le L$,  
the common point
$$
\vec p_{ij} = \vec r_i + s_i\,\vec d_i
$$
is added to the set $\{\vec p_m\}$.

The visualization graph $G = (V,E)$ is then built as
$$
V = \{\vec p_m\},
\qquad
E = \{(\vec p_a,\vec p_b);\ 
\Vert \vec p_a - \vec p_b \Vert < d_{\rm max},\ 
|t_a - t_b| < \Delta t_{\rm max}\} \,.
$$

Here $d_{\rm max}$ is a spatial proximity threshold and $\Delta t_{\rm max}$
is a temporal adjacency threshold.  
The time $t_m$ assigned to point $\vec p_m$
corresponds to the time stamp of the signal segment from which that point was
derived.

Physically, the requirement $\Vert \vec p_a - \vec p_b \Vert < d_{\rm max}$ enforces
that only adjacent parts of a plausible lightning channel can be connected,
and $|t_a - t_b| < \Delta t_{\rm max}$ ensures causal ordering.

The final plotted structure is obtained by drawing all edges
$(\vec p_a,\vec p_b)\in E$ as short line segments:
$$
\Gamma = \bigcup_{(\vec p_a,\vec p_b)\in E}
\{\vec p_a + u(\vec p_b-\vec p_a);\ 0 \le u \le 1\} \,.
\eqmark
$$

\begitems
* It preserves topological continuity of the discharge, avoiding clustering steps typical for point-based mapping \cite[Sibolla2021].
* It suppresses spurious bearings: an isolated noisy line that has no compatible intersections will not generate edges in $E$ and will not appear in the final structure.
\enditems
